\chapter{Conclusions}
\label{ch:conclusions}

\section{Summary}
\label{sec:conclusions-summary}

The goal of this work was to reconstruct a dynamic indoor scene as a time-varying
field of Gaussian primitives that can be rendered interactively from any viewpoint
and at any instant of the captured time. The static variant of this problem is by
now well understood, and the thesis therefore concentrated on the part that is not.
A scene in which a subject moves while the surrounding room stays fixed poses three
coupled difficulties. The reconstruction must represent the same physical surface
across time rather than an unrelated cloud of primitives per frame. It must explain
image changes by motion rather than by repainting, because a purely photometric loss
cannot tell the two apart. And its cost must scale with the complexity of the motion
rather than with the length of the capture, because otherwise every additional
second of video multiplies both the training time and the storage.

The proposed answer rests on one organizing decision, the separation of the scene
into a static background and a dynamic foreground, each reconstructed with the
machinery best suited to it. The background is fitted once by an unmodified
Gaussian Splatting baseline from a single orbiting camera, supplemented by the four
empty-room photographs that the fixed rig records for segmentation, and reaches a
held-out quality of almost $40$\,dB, which confirms that the static half of the
problem needed no new research. The foreground is observed by four fixed, calibrated cameras and
is reconstructed by a custom temporal method built inside the nerfstudio framework.
The method fits a dense cloud on the first frame and then tracks that one cloud
through the remaining frames under two hard constraints. The number of primitives
never changes, and the colour of every primitive stays frozen at its first-frame
value. Under these constraints the only way the optimizer can lower the loss is to
move and reshape the primitives, so the reconstruction is forced to follow the
subject's actual motion. Each frame is warm-started by a constant-velocity
extrapolation of the previous motion, which was measured to cut the steps needed to
reach a fixed quality target by $84$ percent and to lift the converged per-frame
quality by almost six decibels. Primitives that stray from the subject are
reclaimed at every frame boundary by snapping them back onto the visual hull carved
from the camera silhouettes, with a distance margin that reserves this intervention
for genuine runaways.

A deliberate feature of the pipeline is that it contains no learned component
anywhere. The subject is isolated from the room by comparing every frame against a
single photograph of the empty room. The initial study of this idea exposed a
camouflage weakness where the subject visually resembles the background, and the
final version of the segmentation resolves it with a shadow-aware hysteresis rule.
Weak evidence of change is admitted only when it cannot be explained as a cast
shadow and when it is spatially connected to strong evidence. This raised the mask
intersection-over-union from $0.863$ to $0.962$ against ground truth silhouettes
and healed the silhouette holes that previously fragmented the reconstruction in
the early frames. The same silhouettes seed the first frame through a carved visual
hull, and the experiments showed this seed to be a hard requirement rather than a
convenience. A random initialization with the same primitive budget does not merely
converge more slowly, it collapses outright under the masked losses, while matching
the hull's quality by brute force costs several times as many primitives.

Because the tracked representation holds its population and its colours constant,
every remaining attribute becomes a smooth time series, and the thesis exploited
this twice. A temporal smoothness prior in the optimizer damps the frame-to-frame
fluctuation of opacity, scale and rotation by up to two orders of magnitude without
touching the genuine motion of the positions. A frequency-domain compressor then
stores each attribute as a truncated cosine transform along time, which shrinks the
full sequence by a factor of $43$ at matched render fidelity. The final measured
configuration reconstructs the complete $150$-frame capture at full HD resolution
in under eight minutes on a single GPU, reaching the best structural quality of any
configuration evaluated in the thesis, with background leakage more than an order
of magnitude below the earlier baseline. A dedicated web renderer recombines the
per-frame foreground with the static background and plays the result back smoothly,
because switching frames costs only a visibility toggle rather than a data upload.

The evaluation was designed to be honest about its own limits. All object-level
fidelity numbers are bounded by an illumination gap that is a property of the
ground truth rather than of the method, since hiding the room in the reference
renders also removes the light the room bounces onto the subject. The structural
metrics, coverage, leakage and cloud extent therefore carried the discriminating
signal, and every mechanism of the method was validated by producing its failure
mode on demand and measuring the repair.

\section{Possible Further Development}
\label{sec:conclusions-future}

The most consequential next step is the move from synthetic to real captures. The
synthetic environment was chosen deliberately, because exact camera poses make the
measured tracking error attributable to the method alone, and that property made
the controlled experiments of Chapter~\ref{ch:research} possible. A real capture
gives up this exactness. Poses must be estimated, lenses distort, sensors add
noise, and exposure drifts between frames. The pose problem is the deepest of
these, because a pose error and a scene motion produce the same kind of image
change, and the thesis argued that pose noise is silently absorbed into the
estimated motion. A practical route would keep the fixed-camera rig, calibrate it
once with standard photogrammetric targets, and treat the small residual error as a
per-camera bias to be estimated jointly with the first frame. The background
subtraction should transfer to real footage with modest changes, since it already
handles soft shadows and indirect light, but sensor noise will demand a
noise-adaptive threshold and the empty-room plate may need periodic refreshing
under changing daylight.

A second direction concerns the constant-topology assumption itself. Holding the
primitive count fixed is what makes velocity, trajectories and temporal
compression well defined, but it also means the method cannot represent a subject
that enters the scene midway, an object that is picked up and carried, or a person
who leaves. A principled extension would allow topology events at explicit,
detected instants. The silhouette machinery already provides the detection signal,
because a new subject appears as a connected region of change that no existing
primitive explains. Spawning a new, independently tracked cloud for that region
would preserve the per-cloud invariants while lifting the global restriction, and
the compressor would treat each cloud as its own sequence.

The compression study points at a third opportunity. After the smoothness prior,
the positions, opacities and scales compress into a handful of coefficients each,
and the remaining cost is concentrated almost entirely in the rotation channels.
Quaternions resist truncation for a structural reason, since the same physical
rotation has two sign representations and componentwise smoothing ignores the
manifold they live on. Enforcing sign continuity along time, or compressing in a
tangent space around the previous rotation, should release most of that remaining
budget. Beyond rate, the natural target is streaming. The current renderer holds
every frame resident, which is comfortable at the measured sizes but not for
captures that run into the thousands of frames. A windowed loader that decodes the
cosine coefficients on the fly would turn the stored representation directly into
a playable stream.

Several smaller extensions follow the measurements of the results chapter. The
per-frame tracking budget is still a fixed constant, yet the convergence
histograms show that most frames finish long before it is exhausted. Coupling the
early-stopping criterion that already exists in the implementation to the training
schedule would convert the measured headroom into further wall-clock savings
without any quality risk. The constant-velocity motion model could be raised to
second order or replaced by a small temporal filter, which the measured
acceleration residuals suggest would shrink the remaining warm-start error on
fast direction changes. Finally, the illumination gap that limited the evaluation
invites a relightable extension. Because the subject is reconstructed against
black, its appearance currently bakes in the capture lighting. Estimating a
low-order lighting model for the room and dividing it out of the frozen colours
would let the merged scene relight the subject consistently as it moves.

\section{Applications of the Work}
\label{sec:conclusions-applications}

The most direct application of the pipeline is volumetric video. The output of the
method is a sequence of self-contained Gaussian clouds that any modern splatting
viewer can render, and the measured compression brings a full capture down to a
few tens of megabytes, which is well inside the budget of ordinary web delivery.
The included renderer demonstrates the complete playback path in a browser, with
scrubbing and framerate control, and the rendering server can run remotely so that
a thin client only carries the interaction. This makes the system a working
prototype of a free-viewpoint replay tool for telepresence, remote inspection and
immersive media, where a viewer chooses their own camera path through a recorded
performance.

A second family of applications follows from a property that image-based dynamic
reconstruction methods usually lack. Because the primitive set is conserved, the
reconstruction is simultaneously a dense set of point trajectories, and the thesis
showed these trajectories to be coherent enough to read the subject's motion
directly. This opens the door to quantitative motion analysis on top of an
ordinary multi-camera recording. Gait analysis in rehabilitation, technique study
in sport, and ergonomic assessment on a factory floor all need exactly this kind
of dense, markerless trajectory field, and the fixed calibrated rig the method
assumes is a realistic setup for all three environments.

The static-dynamic decomposition also matches the structure of virtual production
and previsualization work in film and interactive media. A set is scanned once at
high quality, performances are captured separately with a small rig, and the two
are recombined freely, with each take of a performance becoming an asset that can
be replayed over the same background from any viewpoint. The same decomposition
serves digital-twin scenarios, where a factory hall or a laboratory is
reconstructed once and moving elements are tracked live against it, and it serves
the simulation side of robotics, where recorded human motion inside a known static
environment provides realistic dynamic obstacles for navigation and manipulation
policies to train against.

Finally, the methodological contributions apply beyond the specific system. The
shadow-aware hysteresis segmentation is a self-contained improvement to classical
background subtraction and is usable wherever a static camera watches a scene with
difficult, camouflaged foregrounds. The visual-hull seeding and the demonstration
that random initialization fails under masked losses are relevant to any
silhouette-supervised splatting pipeline. And the central lesson of the thesis,
that holding topology and appearance fixed turns a hard inverse problem into a
tractable tracking problem while simultaneously making the result compressible,
is a design pattern for dynamic scene representation in general rather than a
detail of this implementation.
